\begin{proposition}[C.3.20 (Proved): the cubic term as an actual relative square]
The four-manifold in C.3.18 can be chosen disjoint from the fixed
sphere. Write it as $M\subset P$, retain its Lefschetz orientation,
and let $\mathcal L$ be the actual principal $\delta$ circle bundle
on $P$, regarded as a complex line bundle. Its restriction to
$\partial M=T_{\gamma uw}$ is trivial. Put
\[
 s_\delta=\langle c_1(\mathcal L|_M,\mathrm{triv})^2,
                         [M,\partial M]\rangle.
\]
This integer is independent of the choice of boundary trivialization.
For the cap of C.3.14, there is one overall sign $\epsilon_{\rm cap}$,
coming from its comparison with the oriented cut-and-reglue trace,
such that
\[
 k_\psi=\epsilon_{\rm cap}\,3s_\delta.
\]
The same relative square, with the compatible Hopf-meridian
orientation, is the linking coefficient of the marked
$\delta=0$ lift of $T_{\gamma uw}$ in $X\setminus F\simeq S^3$.
No numerical value of $s_\delta$ is assumed in this statement.
\end{proposition}
\begin{proof}
First we arrange the claimed disjointness. In the marked product
of C.3.17, the section of the second Tate family avoids $F$.
Consequently the product circle coordinate of the compact sphere
$F$ has image in a compact proper arc of $S^1\setminus\{0\}$.
An orientation-preserving reparametrization of this circle can
compress that arc to length less than $1/12$. Perform this
reparametrization over a neighbourhood of the central first
elliptic fibre, interpolate through circle diffeomorphisms to the
identity, and leave the outer collar fixed. This changes only the
interior circle coordinate and the extending action, not the
actual annular endpoint $\psi$ or its boundary marking.

Rebuild the map of C.3.18 with this coordinate. The twelve
intersection parameters of one Milnor page with the critical
knot are equally spaced in the product circle. Choose the phase
of the page so that all twelve avoid the compressed arc.
Since $F$ lies over the central elliptic fibre, the inverse image
$M$ of this page is now disjoint from $F$. It has the same
boundary, twelve positive nodes, orientation and affine frame
as in C.3.18. Shrink the deleted fixed tube so that $M\subset P$.

Let $g,u,w$ denote the marked circle coordinates on a positive
real regular cusp fibre, let $a$ be the positive argument of
$q$, and let $r=|q|$. Coefficient coordinates have deck rule
\[
 (g,u,w,d,a)\longmapsto(g,u,w-u,d+g,a+1).
\]
Thus the invariant normalized circle connection is
$\alpha=dd-a\,dg$, with Euler curvature
\[
 E=dg\wedge da.
\]
Choose a global connection on $\mathcal L$ agreeing with this
one on the annular collar. This is possible because the displayed
connection is on the actual restricted bundle. Also choose its
curvature near $\partial P$ to be the pullback of the Hopf
curvature on the normal two-sphere. Then $E^2=0$ near that
boundary. Since $H^4(P)=H^3(P)=0$, there is a three-form $B$
with $dB=E^2$, unique up to an exact form. Because
$H^3(\partial P)=0$, it can moreover be chosen zero on a
boundary collar.

The line bundle on the positive real torus is trivialized by
the marked $d=0$ section. On that torus the displayed connection
is zero in this section. Hence relative Chern--Weil theory and
Stokes' theorem give, with $\partial M=-T_{\gamma uw}$,
\[
 s_\delta=\int_M E^2
          =-\int_{T_{\gamma uw}}B.
\]
Changing the trivialization changes the relative Chern class by
the connecting image of $H^1(\partial M;\Z)$. Such a class is
in the radical of the relative intersection pairing: its
pairing with a relative class is the boundary pairing with the
restriction of that class's absolute image, which is zero here.
Products of two connecting classes are zero as well. Thus the
square is independent of the boundary gauge.

We calculate the cut-and-reglue variation of $x^3$. Double $P$
to $B_5=DP$, and extend $\psi$ and all variations by the identity
on the second half. The usual gluing trace gives a bordism of
line-bundle pairs between the difference of the $\psi$-glued
and identity-glued caps and the difference of their mapping
tori. Concretely, cut the gluing collar into two intervals,
take their product with an interval, and exchange the two
opposite pairings across the resulting square; rounding its
corners leaves the two caps and the two mapping-torus boundary
components. The line bundles extend because all four pairings
use the same specified bundle lift. The identity cap is
$S^2\times S^4$, and the identity mapping torus has a line
pulled back from $B_5$; both have zero cubic pairing. Thus the
desired cubic is the mapping-torus variation, up to its one
overall boundary-orientation sign.

On the annulus $\psi:g\mapsto g+\chi(r)$, where $\chi$ decreases
from one to zero. Consequently
\[
 \Delta=\psi^*E-E=d\chi\wedge da.
\]
There is a global one-form $\eta$ with $d\eta=\Delta$ whose
annular value is
\[
 \eta=-(1-\chi)\,da-12\,dg.
\]
Indeed the outer value $-(da+12\,dg)$ extends as a closed form
over the exterior by C.3.16, while the inner value $-12\,dg$
extends over the product cusp by C.3.17. Closed representatives
can be made to agree exactly on collars by subtracting exact
forms. An exact correction makes $\eta$ zero near $\partial P$,
since this boundary has $H^1=0$. This correction does not affect
any integral below, as $B$ vanishes there.

We use the normalized characteristic-form transgression of
\citet[Section 3, Propositions 3.2 and 3.8]{ChernSimons1974E}.
The interpolating curvature on the mapping cylinder is
$E+t\Delta+dt\wedge\eta$. Its endpoint connections can be
matched to the chosen clutching lift by an exact gauge correction,
because $H^1(B_5;\R)=0$. The transgression computation gives
\[
 \int(E+t\Delta+dt\wedge\eta)^3
 =\int_P\eta\wedge(3E^2+3E\Delta+\Delta^2)
 =3\int_P\eta\wedge E^2.
\]
Here $E\Delta=\Delta^2=0$ on the support of $\Delta$. Since
$d(\eta\wedge B)=\Delta\wedge B-\eta\wedge E^2$,
the last expression is $3\int_P\Delta\wedge B$.
Use the orbit orientation
$dg\wedge du\wedge dw\wedge dr\wedge da$, obtained by
taking the positive $\delta$ direction first in the zero-unit
complex model. Integration over the annulus, with $\chi$ going
from one to zero, now gives
\[
 \int_P d\chi\wedge da\wedge B
   =-\int_{T_{\gamma uw}}B=s_\delta.
\]
This proves the asserted factor of three and fixes it without
using a primitive-generator normalization in $H^3(SO)$.

Finally on the total circle bundle over $P$ the form
\[
 \Omega=\alpha\wedge\pi^*E-\pi^*B
\]
is closed. Its integral on the normal Hopf three-sphere is one:
the first term is the normalized Hopf integral and the second
vanishes in the chosen fixed collar. It therefore represents
the integral generator of $H^3(X\setminus F;\Z)$. On the marked
lift of the positive real $T_{\gamma uw}$, the connection and
Euler curvature restrict to zero, and its integral is
$-\int_T B=s_\delta$. This is precisely the claimed linking
coefficient. In particular the square is an invariant of the
actual inclusion, not merely of the abstract genus-one surface.
\end{proof}

\begin{remark}[Scope]
To evaluate this number by a height calculation one must identify
the actual line bundle $\mathcal L|_M$, including its mixed
base--fibre class, with the Poincar\'e line attached to the
original normal-function section. The numerical square of an
auxiliary split or torsor surface does not make that
identification. The signed cancellation with the Pontryagin
term likewise uses the same cut-and-reglue orientation in both
calculations.
\end{remark}
